\chapter{Tensors}\label{ch:tensor}

\laheader{tensor.h, tensor.cc}

\section{Purpose and main idea}\label{sec:tensor:purpose}

\code{Tensor<T>} is a dense tensor of arbitrary rank whose indices are organized in \emph{index groups}.
Each group has its own index range, offset and permutational symmetry (none, symmetric, antisymmetric,
hermitean, antihermitean), and only the permutationally non-equivalent elements of a group are stored.
Indices can optionally carry names (\code{INDEXNAME}) and most operations can address indices by name.
The class provides element access with automatic sign/conjugation handling, contractions implemented via
unwinding to matrices and BLAS \code{gemm}, reordering of indices and groups, (anti)symmetrization by a
general permutation algebra, the HOSVD/Tucker decomposition, conversions to and from \code{NRVec},
\code{NRMat}, \code{NRSMat} and the dense four-index classes of chapter~\ref{ch:fourindex}, and binary,
text and NumPy I/O.

Typical uses are many-body quantities (amplitudes, integrals, density matrices) where antisymmetry
in groups of indices should be exploited for storage.
Limitations stated in the source:
\begin{itemize}
\item the implementation is ``rudimentary''; only \code{Tensor<double>} and \code{Tensor<std::complex<double> >}
	are instantiated (tensor.cc);
\item index types \code{LA_index} and \code{LA_largeindex} are \code{int}, so the total size is limited
	to $2^{31}$ elements (\code{NRVec} limitation);
\item symmetry \emph{between} index groups (e.g.\ $(ij|kl)=(kl|ij)$) is not supported; one would have to
	nest \code{Tensor<Tensor<T> >} or use the four-index classes;
\item covariant and contravariant indices are not distinguished: the header \code{#undef}s
	\code{LA_TENSOR_INDEXPOSITION}, the code guarded by this macro is not active.
\end{itemize}

\section{Storage and concepts}\label{sec:tensor:storage}

\subsection{Index groups}

\begin{lstlisting}
typedef class INDEXGROUP {
public:
  int number;        // number of indices in the group (>0)
  int symmetry;      // -2,-1,0,1,2, see below
  LA_index offset;   // indices run offset ... offset+range-1
  LA_index range;
  bool operator==(const INDEXGROUP &rhs) const; bool operator!=(const INDEXGROUP &rhs) const;
} INDEXGROUP;
\end{lstlisting}
\code{INDEXGROUP} is a plain aggregate without constructor; all members must be set by the user.
If the macro \code{LA_TENSOR_ZERO_OFFSET} is defined before including the header, \code{offset} becomes a
\code{static const} zero (which saves some arithmetic); otherwise any offset, also negative, is allowed.
Index groups are compared memberwise.

\begin{center}
\begin{tabular}{@{}rp{0.3\linewidth}p{0.17\linewidth}p{0.3\linewidth}@{}}
\toprule
\code{symmetry} & meaning & stored elements & odd permutation gives\\
\midrule
0 & no symmetry & $r^m$ & ---\\
1 & symmetric & $\binom{r+m-1}{m}$ & $+x$\\
$-1$ & antisymmetric & $\binom{r}{m}$ & $-x$; zero for a repeated index\\
2 & hermitean (complex \code{T} only) & $\binom{r+m-1}{m}$ & $x^*$\\
$-2$ & antihermitean (complex \code{T} only) & $\binom{r+m-1}{m}$ & $-x^*$\\
\bottomrule
\end{tabular}
\end{center}
Here $m$ is \code{number} and $r$ is \code{range}. For real \code{T} the symmetries $\pm 2$ are illegal
(\code{canonicalize_shape()} raises an error); a group with a single index always gets symmetry 0.
For (anti)hermitean groups the diagonal elements (repeated indices) are stored including their
real/imaginary part which should vanish; this is \emph{not} checked automatically, see
\code{enforce_hermiticity()} and \code{fulfills_hermiticity()}.
The symmetry factors of different groups combine multiplicatively (for complex \code{T} via a small group
multiplication table, so that e.g.\ two conjugations cancel).

\subsection{Storage order}

The element data are kept in the public member \code{NRVec<T> data}. The \emph{leftmost} index is the least
significant, i.e.\ it changes fastest: group 0 is the innermost loop, and inside a non-symmetric group
the first index is the fastest. For a group with $m$ indices $I_0,\dots,I_{m-1}$ and offset subtracted, the
position inside the group is
\begin{align*}
&\textstyle\sum_{k} I_k\, r^k && \text{(no symmetry)},\\
&\textstyle\sum_{k} \binom{I_k+k}{k+1}, \quad I_0\le I_1\le\dots\le I_{m-1} && \text{(symmetric, (anti)hermitean)},\\
&\textstyle\sum_{k} \binom{I_k}{k+1}, \quad I_0< I_1<\dots< I_{m-1} && \text{(antisymmetric)},
\end{align*}
i.e.\ the canonical (stored) element has nondecreasing (increasing) indices; for $m=2$ this is the
usual lower-triangle packing \code{i1*(i1+1)/2+i0} of \code{NRSMat}. The total position is
$\sum_g p_g\,\mathtt{cumsizes}[g]$, where \code{groupsizes[g]} is the number of stored elements of group $g$
and \code{cumsizes[g]} the product of the sizes of groups $0\dots g-1$.
The binomial coefficients are taken from the precomputed table of \code{simplicial()} (miscfunc.h),
which limits symmetric groups to fewer than 16 indices and ranges below 1024.

\subsection{Public data members}

All data members are public:
\begin{lareference}
\lafunc{NRVec<INDEXGROUP> shape}{the index groups; \code{shape.size()==0} for a scalar (rank 0).}
\lafunc{NRVec<T> data}{the stored elements.}
\lafunc{int myrank}{total number of indices (sum of \code{number} over groups), $-1$ for a default-constructed tensor.}
\lafunc{NRVec<LA_largeindex> groupsizes; NRVec<LA_largeindex> cumsizes}{redundant, precomputed from \code{shape} by \code{calcsize()}.}
\lafunc{NRVec<INDEXNAME> names}{index names in flat order (all indices of group 0, then group 1, \dots);
	either empty (unnamed tensor) or of length \code{rank()}.}
\end{lareference}
If \code{shape} is modified directly, \code{calcsize()}, \code{calcrank()} and possibly \code{data.resize()}
must be called by the user to restore consistency.

\subsection{Index types}

\begin{lareference}
\lafunc{typedef int LA_index; typedef int LA_largeindex}{an index value; a position in \code{data}.}
\lafunc{typedef NRVec<LA_index> FLATINDEX}{all indices of an element in flat order.}
\lafunc{typedef NRVec<FLATINDEX> SUPERINDEX}{indices structured by groups: \code{I[g][i]}.}
\lafunc{typedef NRVec<LA_largeindex> GROUPINDEX}{one position (inside the compressed group storage) per group.}
\lafunc{typedef NRMat<LA_index> INDEXMATRIX}{rows are the \code{FLATINDEX}es of all stored elements, in the order of \code{data}.}
\lafunc{struct INDEX {int group; int index;}}{position of one index: group number and index number inside the group (both from 0); constructors \code{INDEX()} and \code{INDEX(int g, int i)}, \code{operator==}.}
\lafunc{typedef NRVec<INDEX> INDEXLIST}{a list of indices.}
\lafunc{struct INDEXNAME {char name[N_INDEXNAME];}}{index name, \code{N_INDEXNAME}=8; constructible from \code{const char *} (copied by \code{strncpy}), so a string literal can be passed wherever an \code{INDEXNAME} is expected; \code{operator=}, \code{==}, \code{!=} compare by \code{strncmp}. Use at most 7 characters so that the name stays NUL-terminated.}
\end{lareference}
\code{LA_traits} are specialized for \code{INDEXNAME}, \code{INDEXGROUP} and \code{INDEX} as plain data,
so that vectors of them can be compared and written/read in binary form (section~\ref{sec:concepts:traits},
\ref{sec:concepts:io}).

\subsection{Reference counting and copy on demand}

\code{Tensor} itself has no reference counter, but all its large members are \code{NRVec}s, so the
scheme of section~\ref{sec:concepts:refcount} applies member by member: the copy constructor and
\code{operator=} share \code{data}, \code{shape}, \code{names} etc.\ (shallow copy), and
\code{copyonwrite()} detaches all of them. There is \emph{no} \code{operator|=}; a deep copy is made by
\code{Tensor<T> b(a); b.copyonwrite();}.
Arithmetic members delegate to \code{NRVec} methods which call \code{copyonwrite()}; contractions call
\code{data.copyonwrite()}; results of restructuring operations are freshly allocated.
An explicit \code{copyonwrite()} is needed before writing through \code{lhs()}, which, like a non-const
\code{operator[]}, raises an error in DEBUG builds when the data are shared. The element reading \code{operator()} is
\code{const} only, so reading a shared tensor is safe. \code{loopover()} and \code{grouploopover()} pass writable pointers
to the callback and therefore detach the data first (one copy if shared, otherwise only a check); use
\code{constloopover()}/\code{constgrouploopover()} for reading only, they never copy. \code{renameindex()},
\code{permute_inside_group()}, \code{enforce_hermiticity()} and \code{flatten()} also detach what they modify.

\section{Construction, assignment, conversion}\label{sec:tensor:construction}

\begin{lareference}
\lafunc{Tensor()}{empty tensor, \code{rank()==-1}.}
\lafunc{explicit Tensor(const T &x)}{scalar (rank 0) with value \code{x}.}
\lafunc{Tensor(const NRVec<INDEXGROUP> &s); Tensor(const INDEXGROUP &g)}{allocate a tensor of the given shape (data not initialized); a single group can be given directly. The shape is canonicalized (\code{canonicalize_shape()}).}
\lafunc{Tensor(const NRVec<INDEXGROUP> &s, const NRVec<INDEXNAME> &newnames); Tensor(const INDEXGROUP &g, const NRVec<INDEXNAME> &newnames)}{the same with index names; \code{newnames} must be empty or of length \code{rank()}.}
\lafunc{Tensor(const NRVec<INDEXGROUP> &s, const NRVec<T> &mydata); Tensor(const NRVec<INDEXGROUP> &s, const NRVec<T> &mydata, const NRVec<INDEXNAME> &newnames)}{use \code{mydata} as the data (shared, not copied); its size must match the shape.}
\lafunc{Tensor(int xrank, const NRVec<INDEXGROUP> &xshape, const NRVec<LA_largeindex> &xgroupsizes, const NRVec<LA_largeindex> xcumsizes, const NRVec<T> &xdata); Tensor(..., const NRVec<INDEXNAME> &xnames)}{low-level constructor setting all members directly without any check.}
\lafunc{Tensor(const Tensor &rhs); Tensor& operator=(const Tensor &rhs)}{shallow copy, see above.}
\lafunc{void resize(const NRVec<INDEXGROUP> &s)}{set a new shape, reallocate data (not initialized) and delete the names.}
\lafunc{void deallocate()}{release data, shape, names and auxiliary arrays.}
\lafunc{void copyonwrite()}{detach all members (deep copy if shared).}
\end{lareference}

\subsection{Conversions}

\begin{lareference}
\lafunc{explicit Tensor(const NRVec<T> &x)}{rank-1 tensor, offset 0, data shared with \code{x}.}
\lafunc{explicit Tensor(const NRMat<T> &x, bool flat=false)}{rank-2 tensor with a copy of the data and offset 0. A square matrix gives one non-symmetric group of 2 indices unless \code{flat} is true; otherwise two groups of one index. Because of the storage order, \code{t(i,j)==x(j,i)}: tensor index 0 is the column index.}
\lafunc{explicit Tensor(const NRSMat<T> &x)}{one group of two indices, data shared with \code{x}; symmetric (symmetry 1) for real, hermitian (symmetry 2) for complex types, as \code{NRSMat} itself (chapter~\ref{ch:smat}); equal to \code{Tensor(NRMat<T>(x))}. If \code{NRSMat} gets the planned \code{hermitian} flag (section~\ref{sec:smat:hermitian}), this constructor has to respect it.}
\lafunc{NRMat<T> matrix(int n=1) const}{reinterpret the data as a matrix whose column index runs over the $n$ leftmost index groups (e.g.\ an index made leftmost by \code{unwind_index()}) and whose row index runs over the remaining groups. The matrix \emph{shares} the data with the tensor. \code{Tensor<T>(x).matrix()} reproduces \code{x}.}
\lafunc{template<fourindexsymtype S, typename T, typename I> Tensor<T> fourindex2tensor(const fourindex_dense<S,T,I> &f)}{conversion from a dense four-index object (chapter~\ref{ch:fourindex}), see the table below; indices get offset 1.}
\lafunc{template<fourindexsymtype S, typename T, typename I> void tensor2fourindex(const Tensor<T> &t, fourindex_dense<S,T,I> &f)}{inverse conversion; shape, ranges, offsets and symmetries are checked.}
\end{lareference}

\code{fourindex2tensor}/\code{tensor2fourindex} are implemented (overloaded) for the following symmetry types.
Since the leftmost tensor index is the least significant, \code{f(i,j,k,l)} corresponds to \code{t(k,l,i,j)}
(for the types marked * this is equal to \code{t(i,j,k,l)} or \code{t(l,k,j,i)} by symmetry):
\begin{center}
\begin{tabular}{@{}p{0.5\linewidth}p{0.46\linewidth}@{}}
\toprule
\code{fourindexsymtype} & tensor shape\\
\midrule
\code{nosymmetry} & 4 groups of 1 index, range $n$\\
\code{twoelectronrealmullikanAB}* & 2 symmetric groups of 2 indices\\
\code{twoelectronrealmullikan}* & 2 symmetric groups of 2 indices (the pair symmetry is lost, data expanded)\\
\code{T2IjAb_aces} & 4 groups of 1 index, ranges \code{nvra, nvrb, noca, nocb}\\
\code{T2ijab_aces} & 2 antisymmetric groups of 2 indices, ranges \code{nvrt}, \code{nocc}\\
\code{antisymtwoelectronrealdiracAB}* & 4 groups of 1 index (data expanded)\\
\code{antisymtwoelectronrealdirac}* & 2 antisymmetric groups of 2 indices (data expanded)\\
\bottomrule
\end{tabular}
\end{center}
In the reverse direction the symmetric matrix of pairs is obtained by symmetrizing the tensor data.

\subsection{NumPy files}

\begin{lareference}
\lafunc{void numpy_write(const char *name, const char *descr="<f8") const}{write a flat tensor (\code{is_flat()}) as a NumPy \code{.npy} file (format version 1.0, C order, the shape is written in reverse so that the NumPy array index order equals the tensor index order). \code{descr} must match \code{T} (\code{"<f8"} for \code{double}, \code{"<c16"} for complex); this is not checked.}
\lafunc{void numpy_read(const char *name)}{read a \code{.npy} file (version 1.0, C or Fortran order) into a flat tensor with offsets 0; the type is not checked.}
\end{lareference}

\section{Element access}\label{sec:tensor:access}

\begin{lareference}
\lafunc{T operator()(LA_index i1...) const}{element by a variable list of \code{rank()} indices (of type \code{int}); returns the value including the symmetry sign/conjugation, 0 for a repeated index in an antisymmetric group.}
\lafunc{T operator()(const FLATINDEX &I) const; T operator()(const SUPERINDEX &I) const}{the same with the indices in a vector.}
\lafunc{Signedpointer<T> lhs(LA_index i1...); Signedpointer<T> lhs(const FLATINDEX &I); Signedpointer<T> lhs(const SUPERINDEX &I)}{element as an l-value: the returned proxy supports \code{=}, \code{+=}, \code{-=} (the value is transformed by the sign/conjugation of the permutation before being stored), \code{*=} and \code{/=} (for an element stored as a conjugate the stored value is multiplied by the conjugated factor, so that the element itself is multiplied by the factor). Writes to elements with a repeated index of an antisymmetric group are ignored (a DEBUG build prints a warning once). Call \code{copyonwrite()} first if the tensor may be shared.}
\lafunc{LA_largeindex index(int *sign, const SUPERINDEX &I) const; LA_largeindex index(int *sign, const FLATINDEX &I) const}{position in \code{data} of the canonical element and the sign: 1, $-1$, 2 (conjugate), $-2$ (negative conjugate), or 0 and position $-1$ for an element which vanishes by antisymmetry. Index ranges are checked in DEBUG builds.}
\lafunc{LA_largeindex vindex(int *sign, LA_index i1, va_list args) const}{the same for a \code{va_list}; ends \code{args}.}
\lafunc{SUPERINDEX inverse_index(LA_largeindex s) const}{canonical indices of the element stored at position \code{s} (not efficient).}
\end{lareference}

\subsection{Looping over all elements}

\begin{lareference}
\lafunc{void loopover(void (*callback)(const SUPERINDEX &, T *), bool skipzeros=false)}{call \code{callback} for every stored element with its canonical indices and a pointer into \code{data}. With \code{skipzeros} the index value 0 is skipped in all indices (useful e.g.\ for indices running $\dots,-2,-1,1,2,\dots$ with a negative offset); the pointer passed to the callback always points to the element of the passed indices (with \code{skipzeros} it is computed from the indices).}
\lafunc{void constloopover(void (*callback)(const SUPERINDEX &, const T *), bool skipzeros=false) const}{the same for reading.}
\lafunc{void grouploopover(void (*callback)(const GROUPINDEX &, T *)); void constgrouploopover(void (*callback)(const GROUPINDEX &, const T *)) const}{loop over all stored elements with the position inside each group's compressed storage instead of the individual indices.}
\lafunc{INDEXMATRIX indexmatrix() const}{matrix whose row $s$ is the flat index of \code{data[s]}; allows a simple \code{for} loop over all elements instead of a callback.}
\end{lareference}

\subsection{Index bookkeeping}

\begin{lareference}
\lafunc{int myflatposition(int group, int index) const; int myflatposition(const INDEX &i) const}{position of an index in the flat index list.}
\lafunc{INDEX myindexposition(int flatindex) const}{inverse of \code{myflatposition}.}
\lafunc{int flatposition(int group, int index, const NRVec<INDEXGROUP> &shape); int flatposition(const INDEX &i, const NRVec<INDEXGROUP> &shape); INDEX indexposition(int flatindex, const NRVec<INDEXGROUP> &shape)}{free versions for a given shape.}
\lafunc{LA_largeindex subindex(int *sign, const INDEXGROUP &g, const NRVec<LA_index> &I)}{position of the index tuple \code{I} inside the storage of one group, with sign as for \code{index()}.}
\lafunc{NRVec<LA_index> inverse_subindex(const INDEXGROUP &g, LA_largeindex s)}{canonical index tuple of position \code{s} in a group.}
\lafunc{FLATINDEX superindex2flat(const SUPERINDEX &I)}{concatenate the groups of a superindex.}
\lafunc{bool zero_in_index(const FLATINDEX &); bool zero_in_index(const SUPERINDEX &); bool zero_in_index(const INDEXMATRIX &, const LA_largeindex row)}{does any index equal 0 (for index ranges with a negative offset where 0 is excluded)?}
\lafunc{bool shares_index(const FLATINDEX &I, const FLATINDEX &J)}{do \code{I} and \code{J} have a common index value?}
\end{lareference}

\section{Arithmetic}\label{sec:tensor:arithmetic}

Binary operations require identical shapes (\code{laerror} otherwise) and, if both tensors are named,
identical names.
\begin{lareference}
\lafunc{Tensor& operator*=(const T &a); Tensor operator*(const T &a) const; Tensor& operator/=(const T &a); Tensor operator/(const T &a) const}{scaling.}
\lafunc{Tensor& operator+=(const Tensor &rhs); Tensor& operator-=(const Tensor &rhs); Tensor operator+(const Tensor &rhs) const; Tensor operator-(const Tensor &rhs) const}{elementwise sum and difference of the stored data.}
\lafunc{Tensor& axpy(const T alpha, const Tensor &rhs)}{\code{*this += alpha*rhs}.}
\lafunc{Tensor operator-() const}{negation; the result has no index names.}
\lafunc{Tensor operator*(const Tensor &rhs) const}{outer product; the indices of \code{rhs} become the less significant (leftmost) ones: \code{(a*b)(J,I)=a(I)*b(J)}. Names are concatenated, so either both or none of the factors must be named.}
\lafunc{T dot(const Tensor &rhs) const}{full contraction $\sum_I \overline{a_I} b_I$ (complex conjugation of \code{*this}); antisymmetric groups are accounted for by the factor $m!$, symmetric and (anti)hermitean groups are not supported (\code{laerror}).}
\lafunc{typename LA_traits<T>::normtype norm() const}{Frobenius norm of the \emph{stored} data (redundant elements are not counted).}
\lafunc{Tensor& conjugateme(); Tensor conjugate() const}{complex conjugation.}
\lafunc{void clear()}{set all elements to zero.}
\lafunc{void randomize(const typename LA_traits<T>::normtype &x)}{random elements, followed by \code{enforce_hermiticity()}.}
\lafunc{void enforce_hermiticity()}{zero the imaginary (hermitean) or real (antihermitean) part of elements with repeated indices in (anti)hermitean groups.}
\lafunc{bool fulfills_hermiticity() const}{check the above condition.}
\end{lareference}

\section{Contractions}\label{sec:tensor:contractions}

Contractions are performed by unwinding the contracted indices of both tensors to the leftmost positions
(expanding them to their full range, \code{unwind_index()}), followed by one matrix multiplication
(\code{cblas_dgemm}/\code{cblas_zgemm}).
The index order of the result is: the remaining indices of the second tensor (\code{rhs}, \code{rhs2}) as
the less significant ones, followed by the remaining indices of the first one (\code{*this} or \code{rhs1});
each remaining group keeps its symmetry (a group from which an index was contracted keeps its symmetry
with one index less). If both tensors are named, the names of the contracted indices must agree and the
result is named accordingly.

In the \code{add...} versions the result is accumulated into \code{*this}:
$\mathtt{*this} \leftarrow \beta\,\mathtt{*this} + \alpha\,\mathrm{contraction}$; with \code{doresize=true}
\code{*this} is first resized to the result shape (this requires \code{beta==0}), otherwise its shape (and
names, if named) must match. \code{conjugate1} conjugates the first tensor, \code{conjugate}/\code{conjugate2}
the second one.

\begin{lareference}
\lafunc{void addcontraction(const Tensor &rhs1, int group, int index, const Tensor &rhs2, int rhsgroup, int rhsindex, T alpha=1, T beta=1, bool doresize=false, bool conjugate1=false, bool conjugate=false)}{contract one index of \code{rhs1} with one index of \code{rhs2}.}
\lafunc{void addcontraction(const Tensor &rhs1, const INDEX &I1, const Tensor &rhs2, const INDEX &I2, ...); void addcontraction(const Tensor &rhs1, const Tensor &rhs2, const INDEXNAME &iname, ...)}{the same with \code{INDEX} positions or with the common index name (same trailing arguments).}
\lafunc{Tensor contraction(int group, int index, const Tensor &rhs, int rhsgroup, int rhsindex, T alpha=1, bool conjugate1=false, bool conjugate=false) const; Tensor contraction(const INDEX &I, const Tensor &rhs, const INDEX &RHSI, ...) const; Tensor contraction(const Tensor &rhs, const INDEXNAME &iname, ...) const}{return $\alpha\,$\code{*this}$\cdot$\code{rhs} contracted over one index.}
\lafunc{void addcontractions(const Tensor &rhs1, const INDEXLIST &il1, const Tensor &rhs2, const INDEXLIST &il2, T alpha=1, T beta=1, bool doresize=false, bool conjugate1=false, bool conjugate2=false); void addcontractions(const Tensor &rhs1, const Tensor &rhs2, const NRVec<INDEXNAME> &names, ...)}{contract the indices \code{il1[i]} with \code{il2[i]} simultaneously in a single matrix multiplication. Finding an optimal contraction order is the responsibility of the user (contracting one index at a time may be cheaper). For empty lists the outer product is added.}
\lafunc{Tensor contractions(const INDEXLIST &il1, const Tensor &rhs2, const INDEXLIST &il2, T alpha=1, bool conjugate1=false, bool conjugate2=false) const; Tensor contractions(const Tensor &rhs2, const NRVec<INDEXNAME> &names, ...) const}{returning versions.}
\lafunc{void addgroupcontraction(const Tensor &rhs1, int group, const Tensor &rhs2, int rhsgroup, T alpha=1, T beta=1, bool doresize=false, bool conjugate1=false, bool conjugate=false); Tensor groupcontraction(int group, const Tensor &rhs, int rhsgroup, T alpha=1, bool conjugate1=false, bool conjugate=false) const}{contract over all indices of a whole group without unwinding them; the groups must have equal shape and be non-symmetric or antisymmetric (for antisymmetric groups the factor $m!$ is included, so the result equals the contraction of the flattened tensors).}
\lafunc{Tensor innercontraction(const INDEXLIST &il1, const INDEXLIST &il2) const; Tensor innercontraction(const INDEX &i1, const INDEX &i2) const; Tensor innercontraction(const NRVec<INDEXNAME> &nl1, const NRVec<INDEXNAME> &nl2) const; Tensor innercontraction(const INDEXNAME &n1, const INDEXNAME &n2) const}{trace(s) inside one tensor: indices \code{il1[i]} and \code{il2[i]} are set equal and summed over; the result is a scalar tensor if no index remains.}
\lafunc{void add_permuted_contractions(const char *antisymmetrizer, const Tensor &rhs1, const Tensor &rhs2, T alpha=1, T beta=1, bool conjugate1=false, bool conjugate2=false)}{for uniquely named tensors: contract all indices whose names occur in both \code{rhs1} and \code{rhs2}, apply the antisymmetrizer given as a string, permute the result to the index order of \code{*this} and accumulate. The string uses the notation of Kucharski and Bartlett: space-separated factors, each a list of index names where \code{/} separates classes and \code{,} separates names within a class; e.g.\ \code{"k/l,m"} is $P(k/lm)=1-P_{kl}-P_{km}$. Indices not mentioned are not permuted; \code{""} means just contraction and reordering. See \code{general_antisymmetrizer} in chapter~\ref{ch:permutation}.}
\end{lareference}

\section{Restructuring}\label{sec:tensor:restructuring}

Permutations are \code{NRPerm<int>} (chapter~\ref{ch:permutation}) and act on group numbers $1\dots n$.
Names are permuted together with the indices unless stated otherwise.

\begin{lareference}
\lafunc{Tensor permute_index_groups(const NRPerm<int> &p) const}{reorder the data permuting whole index groups: group $i$ of the result is group $p_i$ of the source (both counted from 1).}
\lafunc{Tensor permute_index_groups(const NRVec<INDEXNAME> &names) const}{the same with the permutation given by the requested order of the groups' first index names (for a flat tensor: of all names).}
\lafunc{Tensor unwind_index_group(int group) const}{move a group to the leftmost (least significant) position, the others keep their order.}
\lafunc{Tensor unwind_index(int group, int index) const; Tensor unwind_index(const INDEX &I) const; Tensor unwind_index(const INDEXNAME &N) const}{take one index out of its group and make it a separate leftmost non-symmetric group, expanded to its full range; the rest of the group keeps its symmetry. Without data reorganization if the index is already first in a non-symmetric group 0.}
\lafunc{Tensor unwind_indices(const INDEXLIST &il) const; Tensor unwind_indices(const NRVec<INDEXNAME> &names) const}{the same for several indices, which become groups $0,1,\dots$ in the listed order.}
\lafunc{Tensor flatten(int group= -1) const}{expand the given group (all groups for $-1$) into single non-symmetric indices, keeping the flat index order and names.}
\lafunc{Tensor merge_indices(const INDEXLIST &il, int symmetry=0) const; Tensor merge_indices(const NRVec<INDEXNAME> &nl, int symmetry=0) const}{opposite of \code{flatten}: the listed indices (equal range and offset) form a new leftmost group of the given symmetry; for a nonzero symmetry the data are (anti)symmetrized, i.e.\ averaged over the permutations with the appropriate signs.}
\lafunc{Tensor merge_index_groups(const NRVec<int> &groups) const}{permute the listed non-symmetric groups to the front (in the given order, the remaining groups follow in their original order) and merge them into one group; group numbers from 0, repeated numbers raise an error.}
\lafunc{void split_index_group(int group)}{in place, formal (no data movement): split a non-symmetric group into single-index groups.}
\lafunc{void split_index_group1(int group)}{in place, formal: split off the first index of a non-symmetric group.}
\lafunc{void merge_adjacent_index_groups(int groupfrom, int groupto)}{in place, formal: merge consecutive non-symmetric groups of equal range and offset into one.}
\lafunc{void permute_inside_group(int group, const NRPerm<int> &p, bool inverse=false)}{for a group with nonzero symmetry: apply the sign/conjugation of the permutation parity to the data and permute the corresponding names.}
\lafunc{Tensor subtensor1(int i) const}{the slice for the value \code{i} of the last (most significant) index, which must form a group of its own.}
\end{lareference}

\section{Permutation algebra, decompositions and transformations}\label{sec:tensor:other}

\begin{lareference}
\lafunc{void apply_permutation_algebra(const Tensor &rhs, const PermutationAlgebra<int,T> &pa, bool inverse=false, T alpha=1, T beta=0, bool conjugate_by_parity=false)}{$\mathtt{*this} \leftarrow \beta\,\mathtt{*this}$, then for all stored elements $I$ of \code{*this}: $\mathtt{this}(I) \mathrel{+}= \alpha \sum_P c_P\, \mathtt{rhs}(P(I))$, with $P$ acting on the flat index list (inverse permutations if \code{inverse}); with \code{conjugate_by_parity} odd permutations take the complex conjugate. Can represent general (anti)symmetrizers, e.g.\ \code{general_antisymmetrizer} (chapter~\ref{ch:permutation}). \code{*this} may have a symmetry which \code{rhs} lacks. The names of \code{*this} are taken from (or checked against) the names of \code{rhs} permuted by the first permutation. Element-by-element, not optimally efficient.}
\lafunc{void apply_permutation_algebra(const NRVec<Tensor> &rhsvec, const PermutationAlgebra<int,T> &pa, bool inverse=false, T alpha=1, T beta=0)}{the same with \code{rhs} being the outer product of the tensors in \code{rhsvec}, which is not formed explicitly (e.g.\ Grassmann products building a density matrix from cumulants); rather inefficient.}
\lafunc{NRVec<NRMat<T> > Tucker(typename LA_traits<T>::normtype thr=1e-12, bool inverseorder=false)}{HOSVD/Tucker decomposition: \code{*this} is replaced by the flat core tensor and the factor matrices (one per index, columns orthonormal) are returned; singular values below \code{thr} are dropped, which reduces the corresponding range of the core.}
\lafunc{Tensor inverseTucker(const NRVec<NRMat<T> > &x, bool inverseorder=false) const}{rebuild the tensor from the core (\code{*this}) and the matrices returned by \code{Tucker} with the same \code{inverseorder}.}
\lafunc{Tensor linear_transform(const NRVec<NRMat<T> > &x) const}{transform every index of group $g$ by the matrix \code{x[g]} (\code{x[g].ncols()} must equal the range; the new range is \code{x[g].nrows()}), preserving the group symmetries: $t'_{\dots a \dots} = \sum_i x_{ai}\, t_{\dots i\dots}$.}
\lafunc{Tensor linear_transform(const int g, const Tensor<T> &x) const; Tensor linear_transform(const int g, const NRMat<T> &x) const}{transform only group \code{g}, keeping its position and symmetry; \code{x} must be a flat rank-2 tensor (useful for raising/lowering indices) or a matrix.}
\end{lareference}
For the \code{Tensor} argument the contraction runs over its index 0, i.e.\ $t'_{\dots a\dots}=\sum_i x(i,a)\,t_{\dots i\dots}$ in tensor
indexing, which is the same as the matrix convention above for \code{x} built by \code{Tensor<T>(matrix,true)}.

\section{Names and shape queries}\label{sec:tensor:names}

\begin{lareference}
\lafunc{bool is_named() const}{are names set? (\code{laerror} if their number differs from the rank).}
\lafunc{bool is_uniquely_named() const}{named without repetitions.}
\lafunc{void defaultnames(const char *basename="i")}{set names \code{i000}, \code{i001}, \dots}
\lafunc{int findflatindex(const INDEXNAME nam) const}{flat position of the named index, $-1$ if not found (\code{laerror} if unnamed).}
\lafunc{INDEX findindex(const INDEXNAME nam) const; NRVec<INDEX> findindexlist(const NRVec<INDEXNAME> &names) const}{position(s) of named indices; \code{laerror} if not found.}
\lafunc{void renameindex(const INDEXNAME namfrom, const INDEXNAME nameto)}{rename an index (see the copy-on-write remark above).}
\lafunc{int rank() const; LA_largeindex size() const}{number of indices; number of stored elements.}
\lafunc{int calcrank(); LA_largeindex calcsize()}{recompute \code{myrank}, resp.\ \code{groupsizes}, \code{cumsizes} from \code{shape} and return the rank, resp.\ size.}
\lafunc{void canonicalize_shape()}{set symmetry 0 for single-index groups; check the symmetry codes.}
\lafunc{bool is_flat() const; bool is_compressed() const; bool has_symmetry() const; bool has_hermiticity() const}{all groups have one index; some multi-index group has nonzero symmetry; some group has nonzero symmetry; complex \code{T} and some group is (anti)hermitean.}
\end{lareference}

\section{Input and output}\label{sec:tensor:io}

\begin{lareference}
\lafunc{void put(int fd, bool with_names=false) const; void get(int fd, bool with_names=false)}{binary I/O of shape, auxiliary arrays, data and optionally names (section~\ref{sec:concepts:io}); with \code{with_names} the names are always written (an empty list if the tensor has none), so \code{put} and \code{get} must use the same \code{with_names}.}
\lafunc{template<typename T> std::ostream & operator<<(std::ostream &s, const Tensor<T> &x)}{text output: the shape, the names, then one line per stored element with its flat indices and value.}
\lafunc{template<typename T> std::istream & operator>>(std::istream &s, Tensor<T> &x)}{reads the format written by \code{<<}.}
\lafunc{std::ostream & operator<<(std::ostream &s, const INDEXGROUP &x); std::istream & operator>>(std::istream &s, INDEXGROUP &x)}{\code{number symmetry offset range}.}
\lafunc{std::ostream & operator<<(std::ostream &s, const INDEXNAME &x); std::istream & operator>>(std::istream &s, INDEXNAME &x)}{the name as a word.}
\lafunc{std::ostream & operator<<(std::ostream &s, const INDEX &x); std::istream & operator>>(std::istream &s, INDEX &x)}{\code{group index}.}
\end{lareference}

\section{Helpers in the header}

\begin{lareference}
\lafunc{template<typename T> T signeddata(const int sgn, const T &data)}{apply a sign code ($\pm1$, $\pm2$ = with conjugation, 0 gives 0) to a value.}
\lafunc{template<typename T> class Signedpointer}{proxy returned by \code{lhs()}, see section~\ref{sec:tensor:access}.}
\end{lareference}

\section{Examples}\label{sec:tensor:examples}

An antisymmetric rank-3 tensor:
\begin{lstlisting}
int n=4;
INDEXGROUP g;
g.number=3; g.symmetry= -1; g.offset=0; g.range=n;
Tensor<double> t(g);       // stores n(n-1)(n-2)/6 elements
t.clear();
t.lhs(0,1,2) = 1.;         // implies t(1,0,2)=-1 etc.
t.lhs(2,1,3) = 5.;         // stored as t(1,2,3) = -5
const Tensor<double> &tc=t;
cout << tc(1,2,3) << " " << tc(3,2,1) << " " << tc(1,1,2) << endl; // -5 5 0
Tensor<double> tf = t.flatten();   // the same tensor stored in full, n^3 elements
\end{lstlisting}

Contractions of named tensors:
\begin{lstlisting}
INDEXGROUP s; s.number=1; s.symmetry=0; s.offset=0; s.range=n;
NRVec<INDEXGROUP> shape(4); for(int i=0; i<4; ++i) shape[i]=s;
Tensor<double> x(shape), y(shape);
x.randomize(1.); y.randomize(1.);
INDEXNAME xn[4]={"i","j","k","l"}; x.names=xn;
INDEXNAME yn[4]={"m","n","j","i"}; y.names=yn;
INDEXNAME cn[2]={"i","j"};
// z(m,n,k,l) = sum_ij x(i,j,k,l) y(m,n,j,i); remaining indices of y come first
Tensor<double> z = x.contractions(y,NRVec<INDEXNAME>(cn));
// result in a prescribed index order, antisymmetrized in k,l:
// w(k,l,m,n) = z(m,n,k,l) - z(m,n,l,k)
Tensor<double> w(shape);
INDEXNAME wn[4]={"k","l","m","n"}; w.names=wn;
w.add_permuted_contractions("k/l",x,y,1.,0.);
\end{lstlisting}

Tucker decomposition:
\begin{lstlisting}
NRVec<INDEXGROUP> sh(3);
for(int i=0; i<3; ++i) {sh[i].number=1; sh[i].symmetry=0; sh[i].offset=0; sh[i].range=n+i;}
Tensor<double> a(sh); a.randomize(1.);
Tensor<double> a0(a); a0.copyonwrite();       // keep the original
NRVec<NRMat<double> > u = a.Tucker(1e-12);    // a now holds the core tensor
Tensor<double> b = a.inverseTucker(u);        // b == a0
Tensor<double> c = a.linear_transform(u);     // the same
\end{lstlisting}

Conversion from a dense four-index object (chapter~\ref{ch:fourindex}):
\begin{lstlisting}
fourindex_dense<nosymmetry,double,int> f(n);
f.randomize(1.);
Tensor<double> t4 = fourindex2tensor(f);     // f(i,j,k,l) == t4(k,l,i,j), indices from 1
fourindex_dense<nosymmetry,double,int> ff;
tensor2fourindex(t4,ff);
\end{lstlisting}

\section{Pitfalls}\label{sec:tensor:pitfalls}

\begin{itemize}
\item The leftmost index is the fastest; a tensor built from a matrix is its transpose in index notation:
	\code{Tensor<T>(x)(i,j)==x(j,i)}.
\item \code{norm()}, \code{operator+=} etc.\ act on the stored data only; the norm of a compressed tensor
	differs from the norm of its flattened form.
\item Results of contractions have the remaining indices of the \emph{second} tensor first.
\item \code{lhs()} writes require an unshared tensor: call \code{copyonwrite()} after copying.
	There is no \code{operator|=}.
\item Default-constructed \code{INDEXGROUP}s are uninitialized; set all members.
\item Operations based on \code{loopover()} use static global variables to pass data to the callbacks;
	they are neither reentrant nor thread safe.
\item Index names longer than 7 characters are not NUL-terminated.
\item An outer product of a named and an unnamed tensor raises an error (inconsistent number of names).
\end{itemize}
